% Term project report -- CSE8003 Learning Automata (VSA inference)
% Compile from this directory: pdflatex term_proj_report.tex
\documentclass[11pt,a4paper]{article}

\usepackage[T1]{fontenc}
\usepackage[utf8]{inputenc}
\usepackage[margin=1in]{geometry}
\usepackage{amsmath,amssymb}
\usepackage{booktabs}
\usepackage{hyperref}
\usepackage{listings}
\usepackage{xcolor}

\lstset{
  language=Python,
  basicstyle=\ttfamily\small,
  keywordstyle=\color{blue!70!black},
  commentstyle=\color{gray!60},
  stringstyle=\color{green!40!black},
  showstringspaces=false,
  breaklines=true,
  frame=single,
  tabsize=4,
  columns=fullflexible,
}

\title{Variable-Structure Automaton Inference from Interaction Sequences}
\author{Engin Durmaz}
\date{\today}

\begin{document}
\maketitle

\begin{abstract}
This report documents work toward inferring a Variable-Structure Automaton (VSA)
from a unique dataset of state--action--feedback tuples generated in a P-model
environment. We summarize the mathematical background for empirical counts and
conditional estimates, describe the implemented methodology for Tasks~1 and~2
(estimation and optimal-action identification) with references to the Python
codebase, and outline planned experiments for the remaining assignment tasks.
\end{abstract}

\section{Introduction}

The term project asks us to reconstruct, from observed sequences only, the
state--action policy, transition dynamics, penalty structure, and norm of
behavior of an unknown automaton with five states
$\{\phi_1,\ldots,\phi_5\}$, three actions $\{\alpha_1,\alpha_2,\alpha_3\}$,
and binary environment feedback $\beta\in\{0,1\}$ (P-model). Each dataset row is
a tuple
$(\text{sequence\_id},\,\text{step},\,\Phi_{\mathrm{src}}(n),\,\Phi_{\mathrm{dst}}(n),\,\alpha(n),\,\beta(n))$.
This document aligns notation with the course assignment and with the
implementation under \texttt{../proj\_python/}.

\section{Background}

\subsection{Empirical count $N(\phi_i,\alpha_k)$}

The number of times action $\alpha_k$ was taken while the source state was
$\phi_i$ is
\begin{equation}
  N(\phi_i,\alpha_k)
  = \#\bigl\{\, n : \Phi_{\mathrm{src}}(n)=\phi_i \;\wedge\; \alpha(n)=\alpha_k \bigr\}.
\end{equation}
In the CSV, $\Phi_{\mathrm{src}}(n)$ is the column \texttt{phi\_src} and
$\alpha(n)$ is \texttt{alpha}.

\subsection{Policy estimate $P(\alpha_k\mid\phi_i)$}

Empirical policy probabilities normalize the counts in each state:
\begin{equation}
  P(\alpha_k\mid\phi_i)
  = \frac{N(\phi_i,\alpha_k)}{\sum_{k'} N(\phi_i,\alpha_{k'})}
  = \frac{\text{times }\alpha_k\text{ taken in }\phi_i}
         {\text{times }\phi_i\text{ was visited}}.
\end{equation}
The denominator is the total number of actions chosen from state $\phi_i$.

\subsection{Row totals \texttt{row\_totals\_phi}}

In code, for each state $\phi_i$,
\begin{equation}
  \texttt{row\_totals\_phi}[\phi_i] = \sum_k N(\phi_i,\alpha_k),
\end{equation}
i.e.\ the sum of the $N(\phi_i,\cdot)$ row in the count matrix \texttt{N\_phi\_alpha}.

\subsection{Transition and penalty counts (brief)}

Similarly, $N(\phi_i,\alpha_k,\phi_j)$ counts transitions
$\phi_i\to\phi_j$ under action $\alpha_k$, and $N(\alpha_k,\beta)$ counts
(action, feedback) pairs. These yield transition probabilities
$P(\phi_j\mid\phi_i,\alpha_k)$ and penalty probabilities
$c_k = P(\beta=1\mid\alpha_k)$ by row-wise normalization of the appropriate
count tables.

\section{Methodology}

All runnable code lives in \texttt{../proj\_python/}.
\texttt{main.py} orchestrates tasks; each task module can be run standalone
(writing outputs under \texttt{outputs/}) or imported (returning structured
results while still printing to the console).

\subsection{Task 1 --- Estimation}

We load the headerless UTF-8 CSV, build \texttt{N\_phi\_alpha} by grouping on
\texttt{phi\_src} and \texttt{alpha}, build per-action transition count
matrices \texttt{N\_phi\_alpha\_phi}, build \texttt{N\_alpha\_beta}, then
derive \texttt{policy}, \texttt{transition}, and \texttt{penalty} ($c_k$).
The following excerpt is the core link between
Equations~(1)--(3) and the implementation.

\lstinputlisting[firstline=117,lastline=164]{../proj_python/estimation.py}

Consistency checks verify that policy rows and transition rows (where data
exist) sum to unity, and that $c_k$ matches the empirical mean of $\beta$
per action. See \texttt{estimation.py} lines 171--193.

\subsection{Task 2 --- Optimal action identification}

Under the P-model, the optimal action minimizes the penalty probability:
\begin{equation}
  \alpha^\ast \in \arg\min_k c_k,
  \qquad c_k = P(\beta=1\mid\alpha_k).
\end{equation}
The implementation reads $c_k$ from the in-memory \texttt{EstimationResult}
when called from \texttt{main.py}, or from \texttt{penalty.csv} after a
standalone Task~1 run. Ties at the minimum are resolved deterministically
(first label in sorted order).

\lstinputlisting[firstline=110,lastline=131]{../proj_python/optimal_action_identification.py}

\subsection{Tasks 3--6 (planned)}

\begin{itemize}
  \item \textbf{Task 3 ($\varepsilon$-optimality):}
        $\varepsilon = 1 - \min_i P(\alpha^\ast\mid\phi_i)$ using the policy
        matrix from Task~1 and $\alpha^\ast$ from Task~2.
  \item \textbf{Task 4 (simulation):} sample synthetic trajectories from the
        learned policy and transition model; compare marginal distributions to
        the dataset.
  \item \textbf{Task 5 (perturbation):} randomly remove 20\% of transitions
        and re-estimate $c_k$, policy, and $\varepsilon$.
  \item \textbf{Task 6 (model critique):} estimate $c_k$ separately per
        source state and compare across states to assess the P-model assumption.
\end{itemize}
These will extend \texttt{main.py} with new modules mirroring the dual-run
convention used for Tasks~1--2.

\section{Experiments}

\textbf{Data.} The provided file \texttt{../Engin-vsa\_dataset\_full.csv} (one row
per step; Greek state/action labels) is the default input for
\texttt{estimation.py}.

\textbf{Reproducibility.}
\begin{enumerate}
  \item Task~1 only (writes CSVs and \texttt{report.txt}):
        \texttt{cd ../proj\_python \&\& python estimation.py}
  \item Task~2 after Task~1 files exist:
        \texttt{cd ../proj\_python \&\& python optimal\_action\_identification.py}
  \item End-to-end in memory (no intermediate files required):
        \texttt{cd ../proj\_python \&\& python main.py}
\end{enumerate}
Planned experiments for the final report will include the dataset fingerprint
checksum, full numeric tables for Tasks~1--2, and (once implemented) figures
and tables for simulation match, perturbation robustness, and per-state
penalty comparison.

\section{Conclusion}

We fixed notation for empirical counts and policy normalization, and showed
how \texttt{N\_phi\_alpha}, \texttt{row\_totals\_phi}, and \texttt{policy} in
\texttt{estimation.py} implement Equations~(1)--(3). Task~2 selects
$\alpha^\ast$ by minimizing $c_k$ from the same estimation pipeline. The
remaining tasks complete the assignment's data-driven story: optimality
diagnostics, generative validation, robustness under missing transitions, and
critique of the P-model assumption.

\end{document}
